\documentclass{article}
\usepackage{amsmath,amsthm,amssymb}
\usepackage{microtype}

\usepackage{hyperref}
\usepackage{fullpage}


\bibliographystyle{plainurl}% the recommended bibstyle

%\usepackage{subfigure}

%
\usepackage{tikz}
\usetikzlibrary{shapes.symbols,arrows,snakes}
\usetikzlibrary{matrix,calc,intersections,patterns}
\newcommand*\circled[1]{
            \node[shape=circle,draw,inner sep=2pt] (char) {#1};}
\newcommand*\boxd[1]{
            \node[shape=rectangle,draw,inner sep=2pt] (char) {#1};}

\usepackage{enumerate}
\usepackage{algorithm,algorithmic}

\renewcommand{\algorithmicrequire}{\textbf{Initialization:}}
\renewcommand{\algorithmicensure}{\textbf{In round $t \ge 1$ do:}}

\newcommand{\INITIALLY}{\REQUIRE{}}
\newcommand{\ROUND}{\ENSURE{}}

\newcommand{\nop}[1]{}


\newtheorem{thm}{Theorem}
\newtheorem{prop}[thm]{Proposition}
\newtheorem{lem}[thm]{Lemma}
\newtheorem{cor}[thm]{Corollary}
\newtheorem{exo}[thm]{Exercise}
\newtheorem{conj}[thm]{Conjecture}

\DeclareMathOperator{\disc}{disc}
\DeclareMathOperator{\hull}{hull}
\DeclareMathOperator{\affine}{affine}
\DeclareMathOperator{\centroid}{centroid}
\DeclareMathOperator{\vol}{vol}
\DeclareMathOperator{\dec}{dec}
\DeclareMathOperator{\Cube}{Cb}
\DeclareMathOperator{\diam}{diam}
\DeclareMathOperator{\dist}{dist}

\DeclareMathOperator{\Sym}{Sym}
\DeclareMathOperator{\Steiner}{Ste}

\DeclareMathOperator*{\argmin}{arg\,min}
\DeclareMathOperator*{\argmax}{arg\,max}

\newcommand{\N}{\mathcal{N}}

\newcommand{\IR}{\mathbb{R}}
\newcommand{\IN}{\mathbb{N}}
\newcommand{\IZ}{\mathbb{Z}}
\newcommand{\IQ}{\mathbb{Q}}
\newcommand{\IC}{\mathbb{C}}
\newcommand{\M}{\mathcal{M}}
\newcommand{\A}{\mathcal{A}}
\newcommand{\E}{\mathcal{E}}

\newcommand{\Bcnc}{A^{\Delta}_{\mathrm {c}}}                %
\newcommand{\Bunu}{\tilde{B}}            %
\newcommand{\Bms}{\tilde{B}}            %
\newcommand{\Bep}{A^{\Delta}_\mathrm{e}}                %
\newcommand{\Benp}{A^{\Delta}_\mathrm{e}}                %
\newcommand{\Bneone}{A^{\Delta}_\mathrm{r}}                %
\newcommand{\Bnetwo}{A^{\Delta}_\mathrm{r}^\mathrm{HA}}        %
\newcommand{\xb}{\overline{x}}    
\newcommand{\SCC}{\mathcal{C}}
\newcommand{\CP}{\mathcal{P}_\ccirclearrowleft}
%
\newcommand{\IRmax}{\overline{\IR}}
\newcommand{\IRmin}{{\IR}_{\min}}
\newcommand{\wstar}{{w}_{*}}
\newcommand{\ito}{{i\!\to}}
\newcommand{\Pa}{{\mathcal{W}}}

\newcommand{\onevec}{{\mathbf{1}}}

\newcommand{\real}[1]{\mathbf{N}_{\geqslant {#1}}}
\newcommand{\realrem}[2]{\mathbf{N}_{\geqslant {#1}}^{({#2})}}

\renewcommand{\le}{\leqslant}
\renewcommand{\leq}{\leqslant}
\renewcommand{\ge}{\geqslant}
\renewcommand{\geq}{\geqslant}
\newcommand{\one}{\mathbb{1}}


\usepackage{graphicx}


\DeclareMathOperator\In{In}
\DeclareMathOperator\Out{Out}
\DeclareMathOperator\Rcv{Rcv}
\DeclareMathOperator\poly{poly}

%\newtheorem{thm}{Theorem}
%\newtheorem{lem}[thm]{Lemma}
%\newtheorem{cor}[thm]{Corollary}
%\theoremstyle{definition}
%\newtheorem{defn}[thm]{Definition}


\title{Asymptotic Hyperplane Consensus}

\author{Matthias F\"ugger\textsuperscript{1} \and Emmanuel Godard\textsuperscript{2} \and Damien Imbs\textsuperscript{2} \and Thomas Nowak\textsuperscript{3} \and Eloi Perdereau\textsuperscript{2}}

\date{\textsuperscript{1} CNRS, LSV, ENS Paris-Saclay, Universit\'e Paris-Saclay \& Inria\\\textsuperscript{2} Aix Marseille University, CNRS, LIF\\\textsuperscript{3} Universit\'e Paris-Saclay, CNRS}


\begin{document}

\maketitle

\begin{abstract}
We introduce the problem of asymptotic hyperplane consensus, which requires the outputs of processes to converge onto a common hyperplane while remaining inside the convex hull of initial vectors.
This is a relaxation of asymptotic consensus in which outputs have to converge to a single point, i.e., a $0$-dimensional affine subspace.
We show that a large class of algorithms used for asymptotic consensus gracefully degrade to solve asymptotic hyperplane consensus in distributed systems with weaker assumptions on the communication network.
\end{abstract}


\section{Introduction}

We show that the problem of asymptotic consensus to hyperplanes of dimension $d \in \IN$ in network models where
each graph has at most $d$ roots is solvable and give bounds on the convergence rate.

Asymptotic consensus \cite{CBFN15:icalp}

Multidimensional asymptotic consensus \cite{FN18:disc}


\paragraph{Organization.} The rest of the paper is structured as follows.
We introduce the model and the problem in Sections~\ref{sec:model} and~\ref{sec:problem}, respectively.
Section~\ref{sec:algos} describes the class of algorithms that allow to solve asymptotic hyperplane consensus, namely averaging algorithms.
In Section~\ref{sec:analysis} we prove our main result, i.e., that any $\alpha$-secure averaging algorithm solves asymptotic hyperplane consensus with $d$-dimensional vectors in $d$-broadcastable network models.
We give conclusions and perspectives for future work in Section~\ref{sec:conclusion}.

\section{Model}\label{sec:model}

We write $\IN = \{1,2,\dots\}$ and $\IN_0 = \IN \cup \{0\}$.
Further, $[n] = \{1,\dots, n\}$.
For directed graphs $G_1 = ([n],E_1)$ and $G_2 = ([n],E_2)$, we write
  $G_1 \circ G_2$ for the product graph $G = ([n], E)$ with $(i,j) \in E$
  if and only if there exists a $u \in [n]$ such that $(i,u) \in E_1$ and
  $(u,j) \in E_2$.


We assume a synchronous message-passing network with dynamic communication
links, akin to message adversaries~\cite{afek13:asynchrony} or the Heard-Of
model~\cite{charron09:HO}.
That is, every process $i\in[n]$ sends a message in every round $t\in\IN$,
then receives the messages from all incoming neighbors in the (directed)
round-$t$ communication graph $G_t=([n],E_t)$, and updates its local state as a function
of its previous state and the received messages.
Every process receives its own message, i.e., $(i,i)\in G_t$ for all $i\in[n]$.
The execution of a deterministic algorithm is uniquely determined by the
initial states and the sequence of communication graphs.
Although the model is formally synchronous, it has been shown to capture
classical round-based asynchronous~\cite{charron09:HO} and even
Byzantine~\cite{biely07:tolerating,mendes15:multidimensional} models.

A \emph{network model}~$\N$ is a set of communication graphs.
An algorithm \emph{solves a problem} in~$\N$ if it satisifies its specification
in all executions in which all communication graphs are chosen from~$\N$.

A communication graph~$G$ is \emph{$k$-rooted} if there exists a set
$M\subseteq[n]$ of at most~$k$ processes such that every process in~$[n]$
is reachable from~$M$ in~$G$.
In this case, we call~$M$ a \emph{root set} of~$G$.
A network model is $k$-rooted if all its communication graphs are.
Charron-Bost et al.~\cite{CBFN15:icalp} have shown that asymptotic consensus,
i.e., $1$-dimensional asymptotic hyperplane consensus, is
solvable in a network model if and only if it is $1$-rooted.

A communication graph~$G$ is \emph{$k$-broadcastable} if there exists a set
$M\subseteq[n]$ of at most~$k$ processes such that every process in~$[n]$
has an incoming edge from a process in~$M$ in~$G$.
In this case, we call~$M$ a \emph{broadcasting set} of~$G$.
Every $k$-broadcastable communication graph is $k$-rooted, but the converse is
not true.
A network model is $k$-broadcastable if all its communication graphs are.


%TODO: introduce, compare with k = 1

\begin{lem}\label{lem:rooted:to:broadcastable}
Any product of at least~$n^{k+1}$ communication graphs that are $k$-rooted is
$k$-broadcastable.
\end{lem}
\begin{proof}
Let $m \geq n^{k+1}$ and let $G_1,\dots,G_m$ be $k$-rooted.
By the pigeonhole principle, there exists a subsequence $G_1',\dots,G_\ell'$ of at least
$\ell \geq n^{k+1}/\binom{n}{k}\geq n^{k+1}/n^k = n$ communication graphs in which~$M$ is a root set.
By considering the influence sets
\begin{align*}
S_i(t) &= \big\{ j\in[n] \mid (i,j) \in E(G_1'\circ \cdots \circ G_t') \big\} \text{ for every $i\in M$ and}\\
S(t)   &= \bigcup_{i \in M} S_i(t)\enspace,
\end{align*}
it is easy to see that $S(t)$ strictly grows until it covers the set~$[n]$ of all processes.
Thus, since $\lvert S(0)\rvert = k$ and $\ell \geq n-k$, influence set $S(\ell)$ covers~$[n]$.
In particular this means that~$M$ is a broadcasting set in the product graph $G = G_1\circ\cdots\circ G_m$.
\end{proof}

\section{Problem}\label{sec:problem}

For any nonempty finite set $X\subseteq \IR^d$, denote by $\poly(X)$ the
  polyhedron generated by~$X$, i.e., its convex hull.
For any measurable set $A\subseteq \IR^d$, denote its volume in $\IR^d$ by
$\vol(A)$.
We write $\lVert z \rVert = \sqrt{\sum_{k=1}^d z_k^2}$ for the Euclidean norm of any $z\in\IR^d$.
For $A,B\subseteq \IR^d$, denote by
\[
\dist(A,B)
=
\inf \big\{ \lVert a - b \rVert \mid a\in A, b\in B \big\}
\]
the Euclidean distance of~$A$ and~$B$.
By slight abuse of notation we will write $\dist(a,B)$ for $\dist(\{a\},B)$ if
$a\in\IR^d$ and $B\subseteq \IR^d$.
For a set $A\subseteq \IR^d$, we denote by $\bar{A}$ the topological closure
of~$A$.

We say that a sequence of vectors $x(t) \in \IR^d$, $t \in \IN$, \emph{converges onto} 
a set $X \subseteq \IR^d$ if $\displaystyle\lim_{t\to\infty}\dist(x(t),X) = 0$.

In the \emph{$d$-dimensional asymptotic hyperplane consensus problem} every
process~$i$ starts with an initial vector $x_i(0) \in \IR^d$ and outputs a
vector $x_i(t) \in \IR^d$ in every round $t\in\IN$ such that in every execution:
\begin{itemize}
\item There exists an affine hyperplane $E\subseteq \IR^d$ such that all sequences $\big(x_i(t)\big)_t$ converge onto~$E$.
\emph{(Hyperplane Agreement)}
\item All sequences $\big(x_i(t)\big)_t$ converge onto the convex hull of the set of initial vectors.
\emph{(Validity)}
\end{itemize}

\begin{thm}\label{lem:imposs}
The $d$-dimensional hyperplane consensus problem is unsolvable in network
models that are not $d$-rooted.
\end{thm}
\begin{proof}
If a network model is not $d$-rooted, then it contains a communication
graph~$G$ that is not $d$-rooted.
But then~$G$ contains at least $d+1$ strongly connected components without incoming edges.
Let~$M$ be a set of $d+1$ processes from different such components.
Without loss of generality, by permuting the process names, assume that $M =
\{1,\dots,d+1\}$.
For every $i\in M$, let~$A_i$ be the set of processes that reach~$i$ in~$G$.
By choice of~$M$, the sets~$A_i$ are pairwise disjoint.% and $i\in A_i$.

Assume that there exists an algorithm that solves $d$-dimensional hyperplane
consensus in the network model.
We construct an execution of the algorithm as follows.
For every $j\in A_i$, choose the initial vectors $x_j(0) = e_i$ if $1\leq i \leq d$ where~$e_i$
is the $i$\textsuperscript{th} standard basis vector, and $x_j(0) = 0$ if $i = d+1$.
Choose the communication graph of every round~$t$ to be $G_t = G$.
By Validity, for every $i\in M$, since the convex hull of initial values
in~$A_i$ is a singleton and~$i$ does not hear from any process outside
of~$A_i$, the sequence $\big(x_i(t)\big)_t$ converges onto $\{x_i(0)\}$, that is, it
converges to~$x_i(0)$.
But the initial values, and thus the limits, of processes in~$M$ do not
belong to a common hyperplane in~$\IR^d$.
This is a contradiction to Hyperplane Agreement.
\end{proof}


\section{Averaging Algorithms}\label{sec:algos}

An \emph{averaging algorithm} keeps only its current vector~$x_i(t)$ as its
state, which it sends in every round.
After having received the vectors of other processes, it updates its vectors to
a weighted average of the received vectors, see Algorithm~\ref{alg:averaging}.
Different averaging algorithms differ in the choice of weights.
Maybe the most natural averaging algorithm is the \emph{equal neighbor}
algorithm, which assigns the same weight to all received vectors.

\begin{algorithm}[ht]
\begin{algorithmic}[1]
\INITIALLY{}
\STATE $x_i$ is process $i$'s initial vector in~$\IR^d$ 
\ROUND{}
\STATE send $x_i$
\STATE receive vectors $x_j$ for $j\in \In_i(t)$
\STATE determine weights~$a_{ij}$ for received vectors
\STATE 
$\displaystyle x_i \gets \sum_{j\in \In_i(t)} a_{ij} x_j$
\end{algorithmic}
\caption{Averaging algorithm for process $i$}
\label{alg:averaging}
\end{algorithm}

%TODO: introduce In_i(t)

We call an averaging algorithm \emph{$\alpha$-secure} if it guarantees
$a_{ij}\geq \alpha$ in round~$t$ if $j\in \In_i(t)$.
The equal neighbor algorithm is $\alpha$-secure with $\alpha = 1/n$.
  
% safe:
%  and algorithm that gaurantees for each process $i \in \[n]$, that $x_i(t+1)$ is within
%$[ (1-\alpha) m_i(t) + \alpha M_i(t) ,\, (1-\alpha) m_i(t) + \alpha M_i(t) ]$ where
%  $m_i(t)$ is the mimimum and $M_i(t)$ the maximum of the values received by $i$
%  in round $t+1$.

%TODO: refer to ICALP'16 paper, compare with safeness

\section{Analysis}\label{sec:analysis}

We start with some basic notation.
Denote by~$X(t)$ the set
$\{ x_i(t) \mid i \in [n] \}$
of vectors in round~$t$
and 
by
$P(t) = \poly(X(t))$ the generated polyhedron.
It is $\vol(P(t)) = 0$ if and only if~$X(t)$, and thus~$P(t)$, is contained in a hyperplane.
We prove solution of asymptotic hyperplane consensus by showing $\vol(P(t)) \to 0$ as $t\to\infty$.

\begin{thm}\label{thm:poss}
Let $0\leq \alpha < 1$.
In a $d$-broadcastable network model every $\alpha$-secure averaging algorithm solves $d$-dimensional asymptotic hyperplane consensus.
Moreover, we have $\vol(P(T)) \leq \varepsilon$ if $T \geq \alpha^{-d} \log \frac{\vol(P(0))}{\varepsilon}$.
\end{thm}

\begin{cor}
Let $0\leq \alpha < 1$.
In a $d$-rooted network model every $\alpha$-secure averaging algorithm solves $d$-dimensional asymptotic hyperplane consensus.
Moreover, we have $\vol(P(T)) \leq \varepsilon$ if $T \geq n^{d+1} \alpha^{-d} \log \frac{\vol(P(0))}{\varepsilon}$.
\end{cor}

\begin{cor}
In a $d$-broadcastable network model the equal neighbor algorithm solves $d$-dimensional asymptotic hyperplane consensus.
Moreover, we have $\vol(P(T)) \leq \varepsilon$ if $T \geq n^d \log \frac{\vol(P(0))}{\varepsilon}$.
\end{cor}

\begin{cor}
In a $d$-rooted network model the equal neighbor algorithm solves $d$-dimensional asymptotic hyperplane consensus.
Moreover, we have $\vol(P(T)) \leq \varepsilon$ if $T \geq n^{2d+1} \log \frac{\vol(P(0))}{\varepsilon}$.
\end{cor}


The rest of this section is devoted to proving Theorem~\ref{thm:poss}.
%TODO: thus, the following assumptions: alpha-secure algo, maybe others

\subsection{Distance to Hyperplane}

%The distance between a point $z\in \IR^d$ and a hyper-plane
%  $H(v,q) = \{ h\in\IR^d \mid \langle h - q , v \rangle = 0 \}$
%  defined by point $q\in \IR^d$ in the hyper-plane
%  and normal vector $v\in \IR^d\setminus\{0\}$, is given by
%  $\dist(\{z\}, H) = \langle z - q , v\rangle / \lVert v\rVert$.
%We provide the proof for a variant for open half-spaces $H$ for completeness.

The following formula for the distance of a point to a half-space is easy to
prove.

\begin{lem}\label{lem:halfspace:distance:formula}
Let $q\in \IR^d$, $v\in \IR^d\setminus\{0\}$, and
$H = \{ h\in\IR^d \mid \langle h - q , v \rangle < 0 \}$.
Then
$\dist(\{z\}, H) = \langle z - q , v\rangle / \lVert v\rVert$
for all $z \in \IR^d \setminus H$.
\end{lem}
%\begin{proof}
%We show the statement for the variant where $H$ is the
%  hyperplane $\{ h\in\IR^d \mid \langle h - q , v \rangle = 0 \}$.
%The generalization to $H$ being a half-space follows from the observation that
%  the closest point in the (closed) half-space to $z$ lies on the boundary.
%
%Let $z \in \IR^d \setminus H$ and $h_z \in H$ with $h_z \neq z$.
%Let $r = z - h_z \neq 0$.
%Then $\lambda$ in
%\begin{align}
%\lvert z-h_z \rvert = \lambda \frac{r}{\lVert r \rVert}\,.
%\end{align}
%is the distance from $h_z$ to $z$. 
%From $h_z \in \bar{H}$ we have
%\begin{align}
%\langle h_z - q, v \rangle = 0
%\end{align}
%and thus
%\begin{align}
%&\langle z - \frac{\lambda}{\lVert r\rVert} r - q, v \rangle = 0 \Leftrightarrow\\
%&- \frac{\lambda}{\lVert r\rVert} \langle r, v \rangle + \langle z-q, v \rangle = 0 \Leftrightarrow\\
%&\lambda = \langle z-q, v \rangle\frac{\lVert r\rVert}{\langle r, v \rangle}
%\end{align}
%Note that $\lambda(r)$ is minimal for maximal $\frac{\lVert r\rVert}{\langle r, v \rangle}$.
%From Cauchy-Schwarz $\langle r, v \rangle \leq \sqrt{\langle r, r \rangle \langle v, v \rangle}$.
%Thus the above is maximal for $r = v$.
%The lemma follows.
%\end{proof}


For any convex combination algorithm, processes only update their value to a value within the
  polyhedron of received values.
Thus any open half-space $H$ that does not contain points in round $t \ge 0$ will not contain points in
  round $t+1$.
The next lemma shows that the distance of the points to $H$ even increases:
  non-roots increase their distance to $H$ to an amount at least proportional to
  the distance between roots and $H$.

\begin{lem}\label{lem:halfspace:zone}
Let $M$ be a broadcasting set of the communication graph $G(t)=([n],E(t))$ at time $t \ge 1$.
Let $R = \{ x_i(t-1) \mid i\in M\}$ be the set of points of~$M$ at time $t-1$.
Further, let $H \subseteq \IR^d$ be an open half-space such that $X(t-1) \cap H = \emptyset$.
Then $\dist(X(t), H) \geq \alpha \dist(R, H)$.
\end{lem}
\begin{proof}
Let $p \in R$ and $q\in \bar{H}$ such that $\lVert p - q\rVert = \dist(R,\bar{H})$.
Then $\lVert p - q\rVert = \dist(R,H)$.
Because the lemma is trivial if $\dist(R,H)=0$, we assume $p\neq q$ in the rest of the proof.

Let $i\in[n]$.
Then, using $\langle x_j(t-1) - q , p - q\rangle \geq 0$ for all $j\in[n]$ because $X(t-1)\cap H = \emptyset$, we have
\begin{equation}\label{eq:lem:halfspace:zone:ineqs}
\begin{split}
\left\langle x_i(t) - q \, , \, p - q \right\rangle
& =
\left\langle \sum_{j\in [n]} w_{i,j}(t) x_j(t-1) - q \,,\, p - q \right\rangle
=
\left\langle \sum_{j\in [n]} w_{i,j}(t) x_j(t-1) - \sum_{j\in[n]} w_{i,j}(t) q \,,\, p - q \right\rangle
\\ & =
\sum_{j\in [n]} w_{i,j}(t) \left\langle x_j(t-1) - q \,,\, p - q \right\rangle
\geq
\alpha \left\langle x_k(t-1) - q \,,\, p - q \right\rangle
\end{split}
\end{equation}
where $k\in M$ is chosen such that $(k,i)\in E(t)$.

By Lemma~\ref{lem:halfspace:distance:formula},
it is
$\langle r - q , p - q \rangle \geq \lVert p - q\rVert\dist(R,H)$
for all
$r\in R$; in particular for $r = x_k(t-1)$.
But then \eqref{eq:lem:halfspace:zone:ineqs} implies
\begin{equation}\label{eq:lem:halfspace:zone:ineq}
\left\langle x_i(t) - q \, , \, p - q \right\rangle
\geq
\alpha \lVert p - q\rVert \dist(R,H)
\enspace.
\end{equation}
Since $x_i(t)\not\in H$, another application of
Lemma~\ref{lem:halfspace:distance:formula}
transforms~\eqref{eq:lem:halfspace:zone:ineq} into
\begin{equation}
\lVert p - q\rVert \dist(\{x_i(t)\},H)
\geq
\alpha \lVert p - q\rVert \dist(R,H)
\enspace.
\end{equation}
Dividing both sides of the inequality by $\lVert p - q\rVert$ and taking the
minimum over all $i\in [n]$ now concludes the proof.
\end{proof}



\subsection{Symmetric Bodies}

Let $r:\IR \to \IR_+$ be a measurable function.
Then $\Sym(r) \in \IR^d$ is defined as the (full) body obtained by rotating
function $r$ around the first coordinate axis.

Lemma~\ref{lem:volumes} states lower and upper bounds on the volumes of a partitioning
  of such bodies for concave functions $r$.
We start with some technical lemmas.

\begin{lem}\label{lem:concave:line}
Let $a,b,c \in \IR$ with $a<b<c$
and let $r : [a,c]\to \IR$ be a concave function.
Define $f:\IR \to \IR$ to be the unique affine function with $f(b) = r(b)$ and
$f(c) = r(c)$.
Then $f(\xi) \geq r(\xi)$ for all $\xi \in [a,b]$ and $f(\xi) \leq r(\xi)$ for
all $\xi \in [b,c]$.
\end{lem}
\begin{proof}
It is
\begin{equation}
f(\xi)
=
\frac{r(b)-r(c)}{c-b} (c- \xi) + r(c)
\end{equation}
for all $\xi \in \IR$.

By concavity of~$r$, whenever
$\xi \in [a, b]$,
we have
$\alpha = (b - \xi)/(c - \xi) \in [0,1]$ and thus
\begin{equation}
\begin{split}
r(b)
=
r\big((1-\alpha)\xi + \alpha c\big)
\geq
(1-\alpha)r(\xi) + \alpha r(c)
=
\frac{c-b}{c-\xi}\big(r(\xi) - r(c)\big) + r(c)
\enspace,
\end{split}
\end{equation}
which implies
$f(\xi) \geq r(\xi)$ and proves the first part of the lemma.

By concavity of~$r$, 
whenever
$\xi \in [b, c]$,
we have
$\beta = (\xi-b)/(c-b) \in [0,1]$
and thus
\begin{equation}
r(\xi)
=
r\big( (1-\beta)b + \beta c \big)
\geq
(1-\beta) r(b) + \beta r(c)
=
f(\xi)
\enspace,
\end{equation}
which proves the second part of the lemma.
\end{proof}

We will actually only need a weaker linear bound for our purposes:

\begin{lem}\label{lem:concave:line:zero}
Let $a,b,c \in \IR$ with $a<b<c$
and let $r : [a,c]\to \IR$ be a nonnegative concave function.
Define $g:\IR \to \IR$ to be the unique affine function with $g(b) = r(b)$ and
$g(c) = 0$.
Then $g(\xi) \geq r(\xi)$ for all $\xi \in [a,b]$ and $g(\xi) \leq r(\xi)$ for
all $\xi \in [b,c]$.
\end{lem}
\begin{proof}
Using Lemma~\ref{lem:concave:line}, it is sufficient to show that $g(\xi) \geq
f(\xi)$ for all $\xi\leq b$ and $g(\xi) \leq f(\xi)$ for all $\xi\geq b$
where~$f$ is the function defined by
\begin{equation}
f(\xi)
=
\frac{r(b)-r(c)}{c-b} (c- \xi) + r(c)
\enspace.
\end{equation}

We note that
\begin{equation}
g(\xi)
=
\frac{r(b)}{c-b} (c- \xi)
\end{equation}
for all $\xi\in \IR$.

It is $f(b) = g(b)$.
The difference of the slopes of~$f$ and~$g$ is equal to $-r(c)/(c-b)$, which is
nonpositive.
This shows the claimed inequalities.
\end{proof}

This linear bound allows us to show that any $\alpha$-quantile of a symmetric
body on the first axis contains a non-negligible portion of the total volume.

\begin{lem}\label{lem:volumes}
Let $h\in\IR_+$, let $\alpha \in (0,1)$, and let $r:\IR \to \IR_+$ be a measurable function that is concave in $[0,h]$.
Define $r_\mathrm{left}, r_\mathrm{right}:\IR\to \IR_+$ by
$$r_\mathrm{left}(\xi) = \begin{cases}
  r(\xi) & \text{if } \xi \in [0,(1-\alpha)h]\,,\\
  0  & \text{otherwise}
\end{cases}$$
and
$$r_\mathrm{right}(\xi) = \begin{cases}
  r(\xi) & \text{if } \xi \in [(1-\alpha)h,h]\,,\\
  0  & \text{otherwise}
\end{cases}$$
Setting $r_0 = r\big((1-\alpha)h\big)$ and $C_{d-1} = \pi^{(d-1)/2}/\Gamma(\frac{d-1}{2}+1)$,
we have
\begin{equation}
\vol(\Sym(r_\mathrm{left})) = \int_0^{(1-\alpha)h} C_{d-1} r(\xi)^{d-1} \ d\xi
\leq
C_{d-1} \frac{r_0^{d-1}h}{\alpha^{d-1}d} (1 - \alpha^d)
\end{equation}
and
\begin{equation}
\vol(\Sym(r_\mathrm{right})) = \int_{(1-\alpha)h}^{h} C_{d-1} r(\xi)^{d-1} \ d\xi
\geq
C_{d-1}\frac{r_0^{d-1}h}{\alpha^{d-1}d} \alpha^d
\enspace.
\end{equation}
\end{lem}
\begin{proof}
Using the linear function~$g$ from Lemma~\ref{lem:concave:line:zero}, we have
\begin{equation}
\begin{split}
\int_0^{(1-\alpha)h} r(\xi)^{d-1} \ dx
& \leq
\int_0^{(1-\alpha)h} g(\xi)^{d-1} \ dx
=
\int_0^{(1-\alpha)h} \frac{r_0^{d-1}}{\alpha^{d-1}h^{d-1}} (h-\xi)^{d-1} \ d\xi
\\ & =
\frac{r_0^{d-1}}{\alpha^{d-1}h^{d-1}d} \big(h^d - \alpha^d h^d\big)
=
\frac{r_0^{d-1}h}{\alpha^{d-1}d} (1 - \alpha^d)
\end{split}
\end{equation}
and
\begin{equation}
\begin{split}
\int_{(1-\alpha)h}^h r(\xi)^{d-1} \ dx
& \geq
\int_{(1-\alpha)h}^h g(\xi)^{d-1} \ dx
=
\int_{(1-\alpha)h}^h \frac{r_0^{d-1}}{\alpha^{d-1}h^{d-1}} (h-\xi)^{d-1} \ d\xi
\\ & =
\frac{r_0^{d-1}}{\alpha^{d-1}h^{d-1}d} \alpha^d h^d
=
\frac{r_0^{d-1}h}{\alpha^{d-1}d} \alpha^d
\enspace,
\end{split}
\end{equation}
which concludes the proof.
\end{proof}

We can now show contraction of the volume of $P(t)$.


\begin{lem}\label{lem:contract}
%Assume that $\vol(\poly(X(t))) > 0$. 
%Further assume that 
If the communication graph $G(t)$ has a broadcasting set~$M$ with $\lvert M\rvert \leq d$,
then
\begin{equation}
\frac{\vol(P(t))}{\vol(P(t))} \leq 1-\alpha^d
\enspace.
\end{equation}
\end{lem}
\begin{proof}
For brevity, set $P = P(t-1)$ and $P' = P(t)$.

Let $A\subseteq \IR^d$ be a hyperplane that contains all values of processes of the broadcasting set~$M$.
Such a hyperplane exists because $\lvert M\rvert - 1 \leq d - 1$.
Without loss of generality, by rotating and translating the coordinate system accordingly,
we assume that $A = \{z\in\IR^d \mid z_1 = 0\}$.

By performing a Steiner-type symmetrization~\cite{CBFN16:centroid}
of~$P$ along the first axis,
we may also assume that there exists a function $r:\IR\to\IR_+$ with
  $P = \Sym(r)$.
Importantly the transformation preserves the body's volume and convexity,
though it is no longer a polyhedron.
In particular, the function~$r$ is concave on its support.

Define the half-spaces
$A_+ = \{z\in \IR^d \mid z_1 > 0\}$
and
$A_- = \{z\in \IR^d \mid z_1 < 0\}$.
It is sufficient to show
\begin{equation}\label{eq:contract:in:halfspace}
{\vol(P'\cap A_+)}
\leq
(1-\alpha^d)
{\vol(P\cap A_+)}
\end{equation}
for then
\begin{equation}
\vol(P')
=
{\vol(P'\cap A_+)}
+
{\vol(P'\cap A_-)}
\leq
(1-\alpha^d)
\big(
{\vol(P\cap A_+)}
+
{\vol(P\cap A_-)}
\big)
=
(1-\alpha^d)
{\vol(P)}
\end{equation}
by symmetry.

Define $h = \sup\{ z_1 \mid z\in P\}$.
The function~$r$ is concave on the interval~$[0,h]$.
Lemma~\ref{lem:halfspace:zone} applied to the open half-space
$H = \{ z\in \IR^d \mid z_1 > h\}$, we see that
$P' \cap (1-\alpha)H = \emptyset$.
Thus, by the non-expansion property of averaging algorithms, we have
\begin{equation}
\vol(P'\cap A_+) \leq \vol\big((P\cap A_+) \setminus (1-\alpha)H\big) =  \vol(\Sym(r_\mathrm{left}))
\end{equation}
where~$r_\mathrm{left}$ and~$r_\mathrm{right}$ are defined as in Lemma~\ref{lem:volumes}.
We thus conclude
\begin{equation}
\begin{split}
\frac{\vol(P'\cap A_+)}{\vol(P\cap A_+)}
& \leq
\frac{\vol(\Sym(r_\mathrm{left}))}{\vol(\Sym(r_\mathrm{left})) + \vol(\Sym(r_\mathrm{right}))}
=
\frac{1}{1 + \vol(\Sym(r_\mathrm{right}))/ \vol(\Sym(r_\mathrm{left}))}
\\
& \leq
\frac{1}{1 + \alpha^d/(1-\alpha^d)}
=
1 - \alpha^d
\enspace,
\end{split}
\end{equation}
which proves~\eqref{eq:contract:in:halfspace}
and thus the theorem.
\end{proof}

\subsection{Proof of Theorem~\ref{thm:poss}}

By Lemma~\ref{lem:contract}, we have $\vol(P(t)) \to 0$ as $t\to\infty$.
Since the sets~$P(t)$ are nonincreasing for any averaging algorithm, this equivalent to the limit set $P^* = \lim_{t\to\infty} P(t) = \bigcap_{t\geq 0} P(t)$, which is convex, being contained in some affine hyperplane $E\subseteq \IR^d$.
Since $P^* \subseteq P(0)$, it suffices to show $\dist(x_p(t),P^*) \to 0$ for all processes $p\in[n]$ to show both Hyperplane Agreement and Validity.

Suppose by way of contradiction that $\dist(x_p(t),P^*) \not\to 0$ for some process~$p$.
Then there exists some $\varepsilon > 0$ such that $\dist(x_p(t),P^*) \geq \varepsilon$ for infinitely many rounds~$t$.
But then
\begin{equation}
\vol(P(t))
\geq
\vol(\poly(\{x(t)\} \cup P^*))
\geq
\varepsilon \vol_E(P^*) / d
>
0
\end{equation}
for infinitely many~$t$ where $\vol_E(P^*)$ denotes the $(d-1)$-dimensional volume of~$P^*$ as a subset of the hyperplane~$E$.
This is a contradiction to $\vol(P(t)) \to 0$.

The bound on the volume of~$P(t)$ follows by repeated application of the bound in Lemma~\ref{lem:contract}.


\section{Conclusion}\label{sec:conclusion}


\bibliography{agents}


\end{document}

